% Blueprint note: Fixed upstream source versions: \cite{KIPSource5aa845241533,KIPSourcebbac1fc05843}. Interface scope notes are retained as native hgraph node comments; historical migration records are local archives.

% Curated mathematical exposition; compiled dependencies and status are supplied by the publisher.

\section{Sphere Adams differentials}

\begin{theorem}[The one-line differential at h6]
\label{sem:Computation:one-line-differential}
\level{coarse}
\uses{sem:Literature:one-line-differentials,sem:Computation:products-and-powers}
From \cite{Adams60}, as recalled in \cite{LWX126}, Introduction.\par
The designated classes satisfy the nonzero Adams differential
\[d_2(h_6)=h_0h_5^2.\]
% Blueprint note: Derived from the fixed literature one-line theorem on the same sphere.
\lean{KIP126.Computation.Standard.d2_h6}
\end{theorem}

\begin{theorem}[Specified d2 and d3 differentials]
\label{sem:Computation:low-differentials}
\level{coarse}
\uses{sem:Computation:polynomial-values,sem:Computation:products-and-powers}
From \cite{LinProgram}, archived Adams differentials.\par
The following are nonzero differentials on the actual sphere:
\[d_2(h_0^6h_6)=h_0b,\quad d_3(h_4x_{109,12})=h_1x_{122,15,2},\]
\[d_3(h_0^2x_{123,13,2})=h_0^2x_{122,16},\quad d_3(x_{126,4})=h_0^2x_{125,5}.\]
\lean{KIP126.Computation.Standard.d2_h0Six_h6,KIP126.Computation.Standard.d3_h4_x109_12,KIP126.Computation.Standard.d3_h0Square_x123_13_2,KIP126.Computation.Standard.d3_x126_4}
\end{theorem}

\begin{theorem}[A specified d7 differential]
\label{sem:Computation:d7-source}
\level{coarse}
\uses{sem:E2Computation:d7-label,sem:Computation:polynomial-values}
From \cite{LinProgram}, archived Adams differential.\par
For the second-page class $A$ of \ref{sem:E2Computation:d7-label},
\[d_7(A)=h_1x_{121,17}\]
is a nonzero differential in bidegrees $(11,134)\to(18,140)$.
\lean{KIP126.Computation.Standard.d7_source}
\end{theorem}

\begin{theorem}[The two d3 candidates]
\label{sem:Computation:d3-candidates}
\level{coarse}
\uses{sem:E2Computation:d3-labels,sem:Computation:polynomial-values}
From \cite{LinProgram}, archived Adams differential alternatives.\par
For the second-page labels $B,B'$ of \ref{sem:E2Computation:d3-labels},
the nonzero differential on $x_{126,6}$ has one of the two targets
\[d_3(x_{126,6})=h_5x_{94,8}\quad\text{or}\quad d_3(x_{126,6})=h_5x_{94,8}+h_6b.\]

\lean{KIP126.Computation.Standard.d3_x126_6_candidates}
\end{theorem}

\section{Survival, incoming control and higher-page vanishing}

\begin{theorem}[The correction term survives nontrivially]
\label{sem:Computation:correction-survival}
\level{coarse}
\uses{sem:Computation:p-q-correction,sem:SpectralSequence:permanent}
From \cite{LinProgram}, archived Adams calculation.\par
The class $C=e_0\Delta h_6g$ survives to a nonzero element of $E_\infty^{13,137}$.
\lean{KIP126.Computation.Standard.correction_permanent}
\end{theorem}

\begin{theorem}[Incoming differentials to T]
\label{sem:Computation:incoming-control}
\level{coarse}
\uses{sem:Computation:w-t-u,sem:AdamsSynthetic:standard-classes}
If $T$ is hit, its incoming differential is either $d_6(W)=T$ or $d_{12}(h_6^2)=T$, as specified in Fact~\ref{paper:7.6}. Consequently no page $r\geq2$ other than $6$ or $12$ can hit $T$, and no differential from $x_{126,6}$ can have target $T$. The $d_6$ target of $W$ belongs to $\{0,T\}$.
\lean{KIP126.Computation.Standard.no_other_incoming,KIP126.Computation.Standard.table9,KIP126.Computation.Standard.w_d6_targets}
\end{theorem}

\begin{theorem}[Boundary products of the permanent-class candidates]
\label{sem:Computation:product-boundaries}
\level{coarse}
\uses{sem:Computation:p-q-correction,sem:Computation:products-and-powers,sem:SpectralSequence:cycles-boundaries}
From \cite{LinProgram}, archived Adams boundary calculations.\par
The products $Ph_2$ and $Qh_2$ are boundaries by page $2$, in the sense of Definition~\ref{sem:SpectralSequence:cycles-boundaries}.
\lean{KIP126.Computation.Standard.p_h2_boundary,KIP126.Computation.Standard.q_h2_boundary}
\end{theorem}

\begin{theorem}[The required higher-page vanishing components]
\label{sem:Computation:higher-page-vanishing}
\level{coarse}
\uses{sem:AdamsSynthetic:families,sem:Computation:actual-coordinates}
From \cite{LinProgram}, archived Adams page calculations.\par
The actual components $E_5^{11,135}$, $E_4^{12,136}$, $E_4^{12,137}$ and $E_5^{13,138}$ are zero. Here zero means that the corresponding additive module is a subsingleton.
\lean{KIP126.Computation.Standard.e5_stem124_af11,KIP126.Computation.Standard.e4_stem124_af12,KIP126.Computation.Standard.e4_stem125_af12,KIP126.Computation.Standard.e5_stem125_af13}
\end{theorem}

\section{The nu-cofiber and detector differentials}

\begin{theorem}[The cofiber differential and survival window]
\label{sem:Computation:cofiber-results}
\level{coarse}
\uses{sem:Computation:nu-cofiber,sem:Computation:v-x-y,sem:Computation:w-t-u}
From \cite{LinProgram}, archived nu-cofiber Adams calculation.\par
There is a nonzero differential $d_3(A_{C\nu})=i_*Y$ in bidegrees $(8,134)\to(11,136)$. The class $i_*T$ survives to page $6$ and is not hit on any page in the window $2\leq r\leq5$.
\lean{KIP126.Computation.Standard.cofiber_d3,KIP126.Computation.Standard.cofiber_target_to_e6,KIP126.Computation.Standard.cofiber_target_not_hit_through5}
\end{theorem}

\begin{theorem}[The actual detector differential]
\label{sem:Computation:detector-differential}
\level{coarse}
\uses{sem:Computation:detector-comparison,sem:Literature:tmf-d-three}
The detector Adams spectral sequence satisfies
\[d_3(v_2^{16})=B_{19,114}\]
where $v_2^{16}$ and $B_{19,114}$ denote the images of the selected detector classes under the comparison.
% Blueprint note: The legacy Lean name for the target is \texttt{detectorBetaGFour}; the article writes its target as $\beta^5g$. The formal assertion is a differential equation and does not independently assert nonzeroness.
\lean{KIP126.Computation.Standard.detector_d3}
\end{theorem}

\section{Finite Adams records and their certification}

\begin{definition}[Finite Adams claim records]
\label{sem:Computation:table-claim-syntax}
\level{coarse}
\uses{sem:Computation:table-coordinate-syntax,sem:SpectralSequence:differential,sem:SpectralSequence:cycles-boundaries}
From \cite{LinProgram}, archived Adams calculation records.\par
A claim row specifies a spectrum, source and target coordinate lists, a page $r$, and one of four predicates: a differential equation, reaching a finite page, a boundary bound, or the refutation of a differential equation. A valid row has $r\geq2$ and valid coordinates; an equation or refutation also requires $(s\prime,t\prime)=(s+r,t+r-1)$. Validity checks the record's syntax and degrees; the mathematical certification is stated separately.
\lean{KIP126.Computation.Table.Raw.Kind,KIP126.Computation.Table.Raw.Claim,KIP126.Computation.Table.Raw.Claim.valid}
\end{definition}

\begin{theorem}[Mathematical meaning of the certified rows]
\label{sem:Computation:table-certification}
\level{coarse}
\uses{sem:Computation:table-claim-syntax,sem:Computation:table-bases}
From \cite{LinProgram}, archived finite table results.\par
Every row in the fixed claim list decodes successfully and satisfies its stated mathematical predicate: reaching a finite page, being a boundary by a page, an actual differential equation, or the negation of that equation. Each class in the staircase list admits a representative on page $1000$.
% Blueprint note: The staircase records do not assert nonzero permanent survival.
\lean{KIP126.Computation.Table.result,KIP126.Computation.Table.staircase_result}
\end{theorem}



\section{Third-page screening of the second line}
\begin{definition}[The candidates after the first Adams differential]
\label{sem:AdamsCalculations:two-line-d2-candidates}
\level{coarse}
A candidate-list revision of \cite{LWX126}, proof of Corollary 1.5.\par
The calculation uses \cite{Adams60}, Theorem 4.6.1, and \cite{LinExt4}, Theorem 1.3.

An ordered pair $(i,j)$ is a second-line $d_2$ candidate if its standard product is among
\[
 H_{0,2},H_{0,3},H_{2,4},H_{1,j}\ (j\geq3),h_j^2\ (j\geq0),
 H_{2,5},H_{3,5},H_{3,6},H_{0,4}.
\]
\lean{KIP126.AdamsCalculations.TwoLineD2Candidate}
% Blueprint note: Approved source correction (2026-10-18): The first paragraph of the proof of Corollary 1.5 in \cite{LWX126} omits h0 h4 and prints j<=6 already in its E3 list. The corrected candidate predicate retains h0 h4 and every square index. \cite{IWX}, Table 6, p.~222, row (15,2,8), gives the later nonzero d3 on h0 h4. The authors' classical E3 dataset \cite{IWXCharts}, Adams-classical-E3.csv, physical rows 27 and 29 including the header, independently records its nonzero source and target h0 d0. The later HHR exclusion is not an E3 statement. Neither this definition nor the bibliography asserts that the remaining fixed-model screening proof is complete.
\end{definition}

\begin{theorem}[The Adams--Lin--Leibniz candidate screening]
\label{sem:AdamsCalculations:two-line-d2-screening}
\level{coarse}
A candidate-screening revision of \cite{Adams60}, Theorem 4.6.1; \cite{LinExt4}, Theorem 1.3.\par

If $i\leq j$, $H_{i,j}\ne0$ on $E_2$, and $H_{i,j}$ has a third-page continuation, then $(i,j)$ is one of the candidates above.
\lean{KIP126.AdamsCalculations.two_line_d2_candidates}
% Blueprint note: Detailed original sources: \cite{Adams58}, Theorem 2.2(iii), p.~182, supplies multiplicative Leibniz. \cite{Adams60}, Theorem 4.6.1, pp.~99--100, supplies the secondary-operation relation whose Adams interpretation is d2(h_j)=h0 h_(j-1)^2 for j>=4; the modern d2 formula is an adaptation, not a verbatim statement of that theorem. \cite{LinExt4}, Theorems 1.2--1.3, pp.~460--461, supplies the products and relations on the first four lines. These are the primary ingredients for a derived calculation, not a literal theorem of Lin.
% Blueprint note: Independent source audit found (0,4) to be a counterexample to the earlier candidate list. The product is nonzero by \cite{Adams60}, Theorem 2.5.1(2), p.~45. Lin Theorem 1.3(1), at i=0, gives h0^2 h3^2=0, hence Leibniz gives d2(h0 h4)=0. Its later nonzero d3 is explicit in \cite{IWX}, Table 6, p.~222. The user approved the correction: add (0,4), supply its separate nonzero d3 Literature interface, and add its exclusion branch in Main. The final permanent-survivor statement is unchanged, and the one additional upstream responsibility raises the total from 309 to 310. The corrected screening and source-to-fixed-model comparison remain a registered upstream proof obligation; citation verification and finite checks do not prove the unbounded interface.
\end{theorem}
